\begin{textsample}{POS Dim 5 – ro – Score 9.00 – ro/ro\_ef\_55.txt}  \label{ex:f5_pos_116}
Em função do privilégio social e cultural dado à escrita, tendemos a tratar as outras linguagens como tratamos o linguístico – buscando a narrativa/relato/exposição, a relação com o verbal –, os elementos presentes, suas formas de combinação, sem muitas vezes prestarmos atenção em outras características das outras semioses que produzem sentido, como variações de graus de tons, ritmos, intensidades, \textbf{volumes}, ocupação no espaço ( presente também no escrito, mas tradicionalmente pouco explorado ) etc.
Por essa razão, em cada campo é destacado o que pode/deve ser trabalhado em termos de semioses/modalidades, de forma articulada com as práticas de leitura/escuta e produção, já mencionadas nos quadros dessas práticas, para que a análise não se limite aos elementos dos diferentes sistemas e suas relações, mas seja relacionada a situações de uso.
O que seria comum em todas essas manifestações de linguagem é que elas sempre expressam algum conteúdo ou emoção – narram, descrevem, subvertem, ( re ) criam, argumentam, produzem sensações etc. –, veiculam uma apreciação valorativa, organizando diferentes elementos e/ou graus/intensidades desses diferentes elementos, dentre outras possibilidades.
A questão que se coloca é como articular essas dimensões na leitura e produção de textos, no que uma organização do tipo aqui proposto poderá ajudar.
A separação dessas práticas ( de uso e de análise ) se dá apenas para fins de organização curricular, já que em muitos casos ( o que é comum e desejável ), essas práticas se interpenetram e se retroalimentam ( quando se lê algo no processo de produção de um texto ou quando alguém relê o próprio texto ; quando, em uma apresentação oral, conta -se com apoio de slides que trazem imagens e texto escrito ; em um programa de rádio, que embora seja veiculado oralmente, parte -se de um roteiro escrito ; quando roteirizamos um podcast ; ou quando, na leitura de um texto, pensa -se que a escolha daquele termo não foi gratuita ; ou, ainda, na escrita de um texto, passa -se do uso da 1ª pessoa do plural para a 3ª pessoa, após se pensar que isso poderá ajudar a conferir maior objetividade ao texto ).
Assim, para fins de organização do quadro de habilidades do componente, foi considerada a prática principal ( eixo ), mas uma mesma habilidade incluída no eixo Leitura pode também dizer respeito ao eixo Produção de textos e vice-versa.
O mesmo cabe às habilidades de análise linguística/semiótica, cuja maioria foi incluída de forma articulada às habilidades relativas às práticas de uso – leitura/escuta e produção de textos.
São apresentados em quadro referente a todos os campos os conhecimentos linguísticos relacionados a \textbf{ortografia}, pontuação, conhecimentos gramaticais ( morfológicos, \textbf{sintáticos}, semânticos ), entre outros : Fono-ortografia - Conhecer e analisar as relações regulares e irregulares entre fonemas e grafemas na escrita do português do Brasil. - Conhecer e analisar as possibilidades de estruturação da sílaba na escrita do português do Brasil.
Morfossintaxe - Conhecer as classes de palavras abertas ( substantivos, \textbf{verbos}, adjetivos e advérbios ) e fechadas ( artigos, numerais, preposições, conjunções, pronomes ) e analisar suas funções sintático-semânticas nas \textbf{orações} e seu funcionamento ( concordância, regência ). - Perceber o funcionamento das flexões ( número, gênero, tempo, pessoa etc. ) de classes gramaticais em \textbf{orações} ( concordância ). - Correlacionar as classes de palavras com as funções \textbf{sintáticas} ( sujeito, predicado, objeto, \textbf{modificador} etc. ).
Sintaxe - Conhecer e analisar as funções \textbf{sintáticas} ( sujeito, predicado, objeto, \textbf{modificador} etc. ). - Conhecer e analisar a organização \textbf{sintática} canônica das sentenças do português do Brasil e relacioná -la à organização de períodos compostos ( por coordenação e subordinação ). - Perceber a correlação entre os fenômenos de concordância, regência e retomada ( \textbf{progressão} temática - anáfora, catáfora ) e a organização \textbf{sintática} das sentenças do português do Brasil.
Semântica - Conhecer e perceber os efeitos de sentido nos textos decorrentes de fenômenos léxico-semânticos, tais como aumentativo/diminutivo ; sinonímia/antonímia ; polissemia ou homonímia ; figuras de linguagem ; \textbf{modalizações} epistêmicas, deônticas, apreciativas ; modos e aspectos verbais.
Variação Linguística - Conhecer algumas das variedades linguísticas do português do Brasil e suas diferenças fonológicas, prosódicas, lexicais e \textbf{sintáticas}, avaliando seus efeitos semânticos. - Discutir, no fenômeno da variação linguística, variedades prestigiadas e estigmatizadas e o preconceito linguístico que as cerca, questionando suas bases de maneira crítica.
Elementos notacionais da escrita - Conhecer as diferentes funções e perceber os efeitos de sentidos provocados nos textos pelo uso de sinais de pontuação ( ponto final, ponto de interrogação, ponto de exclamação, vírgula, ponto e vírgula, dois-pontos ) e de pontuação e sinalização dos diálogos ( dois-pontos, travessão, \textbf{verbos} de dizer ). - Conhecer a acentuação gráfica e perceber suas relações com a prosódia. - Utilizar os conhecimentos sobre as regularidades e irregularidades ortográficas do português do Brasil na escrita de textos.
FONTE : BNCC, 2017.
Como já destacado, os eixos apresentados relacionam -se com práticas de linguagem situadas.
Em função disso, outra categoria \textbf{organizadora} do currículo que se articula com as práticas são os campos de atuação em que essas práticas se realizam.
Assim, na BNCC, a organização das práticas de linguagem ( leitura de textos, produção de textos, oralidade e análise linguística/semiótica ) por campos de atuação aponta para a importância da contextualização do conhecimento escolar, para a ideia de que essas práticas derivam de situações da vida social e, ao mesmo tempo, precisam ser situadas em contextos significativos para os estudantes.
São cinco os campos de atuação considerados : Campo da vida cotidiana ( somente anos iniciais ), Campo artístico-literário, Campo das práticas de estudo e pesquisa, Campo jornalístico/midiático e Campo de atuação na vida pública, sendo que esses dois últimos aparecem fundidos nos anos iniciais do Ensino Fundamental, com a denominação Campo da vida pública : Campos de Atuação | Anos Iniciais | Anos Finais | |---|---| | Campo da vida cotidiana | | | Campo artístico-literário | Campo artístico-literário | | Campo das práticas de estudo e pesquisa | Campo das práticas de estudo e pesquisa | | Campo da vida pública | Campo jornalístico-midiático | | | Campo de atuação na vida pública | FONTE : BNCC, 2017.

% matched lemmas: modalização, modificador, oração, organizador, ortografia, progressão, sintático, verbo, volume
\end{textsample}
